% Part: first-order-logic
% Chapter: proof-systems
% Section: introduction

\documentclass[../../../include/open-logic-section]{subfiles}

\begin{document}

\iftag{FOL}
      {\olfileid{fol}{prf}{int}}
      {\olfileid{pl}{prf}{int}}

\olsection{مقدمة}

للأنساق المنطقية عادةً دلالة ونظام !!a{derivation} معًا. وتتناول
الدلالة مفاهيم مثل الصدق، وقابلية الإرضاء، والصلاحية، والاستلزام. أما
الغرض من أنظمة ال!!{derivation} فهو توفير طريقة تركيبية محضة لإثبات
الاستلزام والصلاحية. وهي تركيبية محضة بمعنى أن ال!!{derivation} في
مثل هذا النظام كائن تركيبي منتهٍ، يكون عادةً متتاليةً (أو ترتيبًا
منتهيًا آخر) من ال!!{sentence}s أو ال!!{formula}s. وتمتاز أنظمة
ال!!{derivation} الجيدة بإمكان التحقق آليًا من «صحة» أي متتالية أو
ترتيب معطى من ال!!{sentence}s أو ال!!{formula}s.

كانت أبسط أنظمة ال!!{derivation} في منطق الرتبة الأولى (وأولها
تاريخيًا) أنظمةً \emph{بديهية}. وتُعدّ متتالية من ال!!{formula}s
!!a{derivation} في مثل هذا النظام إذا كانت كل !!{formula} مفردة فيها
إما من مجموعة ثابتة من «البديهيات»، وإما ناتجة من !!{formula}s تسبقها
في المتتالية بتطبيق واحدة من مجموعة ثابتة من «قواعد الاستدلال»؛ ويمكن
التحقق آليًا مما إذا كانت !!a{formula} بديهية، ومما إذا كانت ناتجةً
على الوجه الصحيح من !!{formula}s أخرى بإحدى قواعد الاستدلال. ويسهل
وصف أنظمة ال!!{derivation} البديهية---كما يسهل تناولها على مستوى ما
وراء النظرية---لكن قراءة ال!!{derivation}s فيها وفهمها وإنتاجها أمور
صعبة.

وقد طُوّرت أنظمة !!{derivation} أخرى بغرض تيسير بناء
ال!!{derivation}s أو فهم ال!!{derivation}s بعد اكتمالها. ومن أمثلتها الاستنتاج
الطبيعي، وأشجار الصدق المعروفة أيضًا ببراهين الجداول الدلالية، وحساب
التتابعيات. وصُممت بعض أنظمة ال!!{derivation} خصوصًا لأغراض المكننة؛
فطريقة الحَلّ، مثلًا، يسهل تنفيذها برمجيًا (لكن فهم
!!{derivation}sها يكاد يكون مستحيلًا). وتمثل معظم أنظمة ال!!{derivation}
الأخرى ال!!{derivation}s في صورة أشجار من ال!!{formula}s بدلًا من
المتتاليات. وهذا ييسّر رؤية أجزاء !!a{derivation} التي يعتمد بعضها
على بعض.

وهكذا تقدم أنظمة ال!!{derivation} المختلفة، لمنطق معطى كمنطق الرتبة
الأولى، شروحًا مختلفة لما يعنيه أن تكون !!a{sentence}
\emph{مبرهنةً}، ولما يعنيه أن تكون !!a{sentence} !!{derivable} من
جمل أخرى. وأيًا تكن الطريقة (بواسطة !!{derivation}s بديهية، أو
استنتاجات طبيعية، أو !!{derivation}s تتابعية، أو أشجار صدق، أو
تفنيدات بطريقة الحَلّ)، فنحن نريد لهذه العلاقات أن تطابق المفهومين
الدلاليين للصلاحية والاستلزام. ولنكتب $\Proves !A$ بمعنى «$!A$ مبرهنة»،
و«$\Gamma \Proves !A$» بمعنى «$!A$ !!{derivable} من~$\Gamma$».
وأيًا يكن تعريف~$\Proves$، نريده أن يطابق~$\Entails$، أي:
\begin{enumerate}
\item $\Proves !A$ إذا وفقط إذا كان $\Entails !A$
\item $\Gamma \Proves !A$ إذا وفقط إذا كان $\Gamma \Entails !A$
\end{enumerate}
يسمى اتجاه «فقط إذا» أعلاه \emph{السلامة}. ويكون نظام
!!^a{derivation} سليمًا إذا كانت !!{derivability} تضمن الاستلزام (أو
الصلاحية). ولا بد لكل نظام !!{derivation} جدير بالاعتبار أن يكون
سليمًا؛ فلا فائدة إطلاقًا من أنظمة ال!!{derivation} غير السليمة. ذلك
أن الغرض كله من !!a{derivation} هو تقديم ضمان تركيبي للصلاحية أو
الاستلزام. وسنثبت السلامة لأنظمة ال!!{derivation} التي نعرضها.

أما الاتجاه العكسي، «إذا»، فمهم أيضًا، ويسمى \emph{الاكتمال}. يكون
نظام ال!!{derivation} الكامل قويًا بما يكفي لإظهار أن $!A$ مبرهنة
كلما كانت $!A$ صالحة، وأن $\Gamma \Proves !A$ كلما كان
$\Gamma \Entails !A$. وإثبات الاكتمال أصعب، ولا توجد لبعض الأنساق
المنطقية أنظمة !!{derivation} كاملة. أما منطق الرتبة الأولى فله مثل
هذه الأنظمة. وكان كورت غودل أول من أثبت، في أطروحته سنة 1929، اكتمال
نظام !!a{derivation} لمنطق الرتبة الأولى.

ومن المفاهيم الأخرى المرتبطة بأنظمة ال!!{derivation} مفهوم
\emph{الاتساق}. تسمى مجموعة من ال!!{sentence}s غير متسقة إذا أمكن
!!{derive} أي شيء منها، ومتسقة فيما عدا ذلك. وعدم الاتساق هو المقابل
التركيبي لعدم قابلية الإرضاء: فمثلما لا تصلح المجموعات غير القابلة
للإرضاء نظرياتٍ جيدة، لا تصلح لذلك أيضًا المجموعات غير المتسقة من
ال!!{sentence}s، إذ يشوبها خلل أساسي. وقد لا تكون المجموعات المتسقة
من ال!!{sentence}s صادقةً أو نافعة، لكنها على الأقل تتجاوز ذلك الحد
الأدنى من النفع المنطقي. وقد يختلف، باختلاف أنظمة ال!!{derivation}،
التعريف الدقيق لاتساق مجموعات ال!!{sentence}s؛ لكننا نريد للاتساق،
مثل~$\Proves$، أن يطابق مقابله الدلالي، وهو قابلية الإرضاء. ونريد أن
تكون~$\Gamma$ متسقة إذا وفقط إذا كانت قابلة للإرضاء. وهنا يعادل اتجاه
«فقط إذا» الاكتمال (إذ يضمن الاتساق قابلية الإرضاء)، ويعادل اتجاه
«إذا» السلامة (إذ تضمن قابلية الإرضاء الاتساق). والواقع أن صيغتي
السلامة والاكتمال متكافئتان في منطق الرتبة الأولى الكلاسيكي.

\end{document}
